\documentclass[10pt]{mypackage}

\usepackage[uprightgreeks,light]{kpfonts}
\renewcommand{\coloneq}{\coloneqq}
%\usepackage{homework}
\usepackage{notes}

\usepackage[ backend=bibtex, style = alphabetic, sorting=ynt ]{biblatex}
\addbibresource{all_references.bib}

\usepackage{parskip}

\fancyhf{}
\fancyhead[R]{Avinash Iyer}
\fancyhead[L]{The Standard Representation}
\fancyfoot[C]{\thepage}

\setcounter{secnumdepth}{0}

\begin{document}
\RaggedRight
\section{The Standard Representation}%
Before we may start understanding the von Neumann algebra of countable Borel equivalence relations, we discuss the standard representation of a finite von Neumann algebra $\left( M,\tau \right)$.
\begin{proposition}
  Let $M$ be a finite von Neumann algebra with trace $\tau$. Then, the map $x\mapsto x^{\ast}$ extends to an anti-linear isometry $J\colon L_2\left( M,\tau \right)\rightarrow L_2\left( M,\tau \right)$ with $JMJ = M'\cap B\left( L_2\left( M,\tau \right) \right)$. Here, $L_2\left( M,\tau \right)$ refers to the GNS representation of $M$ with respect to the faithful state $\tau$.
\end{proposition}
\begin{proof}
  For any $x\in M$, we have
  \begin{align*}
    \norm{x^{\ast}}_{\tau}^2 &= \tau\left( xx^{\ast} \right)\\
                             &= \tau\left( x^{\ast}x \right)\\
                             &= \norm{x}_{\tau}^2.
  \end{align*}
  Thus, this extends to an anti-linear isometry $J\colon L_2\left( M,\tau \right)\rightarrow L_2\left( M,\tau \right)$. We have $ \iprod{J\xi}{J\eta} = \iprod{\xi}{\eta} $ for all $\xi,\eta\in L_2\left( M,\tau \right)$.

  For any $\xi\in L_2\left( M,\tau \right)$, define the unbounded operators $L_{\xi}^{0}\colon M1_{\tau}\rightarrow L_2\left( M,\tau \right)$ and $R_{\xi}^{0}\colon M1_{\tau}\rightarrow L_2\left( M,\tau \right)$ by
  \begin{align*}
    L_{\xi}^0\left( x1_{\tau} \right) &= \left( Jx^{\ast}J \right)\xi\\
    R_{\xi}^{0}\left( x1_{\tau} \right) &= x\xi.
  \end{align*}
  Here, $1_{\tau}$ refers to the cyclic vector for the representation $\tau$.

  If $\left( x_i \right)_i\in M$ and $\eta\in L_2\left( M,\tau \right)$ are such that $\norm{x_i}_{\tau}\rightarrow 0$ and $L_{\xi}^0\left( x_i1_{\tau} \right)\rightarrow\eta$, then for all $z\in M$, we have
  \begin{align*}
    \left\vert \iprod{\eta}{z} \right\vert &= \lim_{i} \left\vert \iprod{\left( Jx_i^{\ast}J \right)\xi}{z} \right\vert\\
                                           &= \lim_{i} \left\vert \iprod{\xi}{zx_i^{\ast}} \right\vert\\
                                           &\leq \lim_{i} \norm{\xi}_{\tau}\norm{z}_{\op}\norm{x_i}_{\tau}\\
                                           &= 0,
  \end{align*}
  meaning that $L_{\xi}^0$ is closable. Completing the graph gives rise to a closed operator $L_{\xi}$, and similarly $R_{\xi}^0$ is closable with corresponding operator $R_{\xi}$. We have that $L_{x1_{\tau}} = x$ for all $x\in M$.

  Now, if $x,y\in M$ and $\xi\in L_2\left( M,\tau \right)$, then
  \begin{align*}
    \iprod{R_{J\xi}\left( x1_{\tau} \right)}{y1_{\tau}} &= \iprod{xJ\xi}{y1_{\tau}}\\
                                                        &= \iprod{y^{\ast}x1_{\tau}}{\xi}\\
                                                        &= \iprod{x1_{\tau}}{R_{\xi}\left( y1_{\tau} \right)},
  \end{align*}
  meaning that $R_{J\xi} = R_{\xi}^{\ast}$. Additionally,
  \begin{align*}
    \left( JL_{\xi}J \right)\left( x1_{\tau} \right) &= xJ\xi\\
                                                     &= R_{J\xi} \left( x1_{\tau} \right)\\
                                                     &= R_{\xi}^{\ast}\left( x1_{\tau} \right),
  \end{align*}
  so that $JL_{\xi}J = R_{\xi}^{\ast}$.

  If we set $ \widetilde{M} = \set{L_{\xi} | L_{\xi}\in B\left( L_2\left( M,\tau \right) \right)} $ and $ \widetilde{N} = \set{R_{\xi} | R_{\xi}\in B\left( L_2\left( M,\tau \right) \right)} $, then we have $J \widetilde{M}J = \widetilde{N}$.

  Therefore, it suffices to show that $M = \widetilde{M}$ and $M' = \widetilde{N}$. For this, we observe that $M\subseteq \widetilde{M}$ and $ \widetilde{N}\subseteq M' $. For any $x\in M'$, let $\xi = x1_{\tau}$. For any $y\in M$, we then have
  \begin{align*}
    R_{\xi}\left( y1_{\tau} \right) &= y\xi\\
                                    &= yx1_{\tau}\\
                                    &= xy1_{\tau},
  \end{align*}
  meaning that $R_{\xi} = x$, so $M'\subseteq \widetilde{N}$. For any $L_{\xi}\in \widetilde{M}$, we then have $L_{\xi}\in \widetilde{N}' = M'' = M$, so $\widetilde{M}\subseteq M$.
\end{proof}
We observe that by considering the maps $R_{\xi}$ and $L_{\xi}$ as left and right multiplication by $M$, the standard representation can be considered as a $M$-$M$-bimodule. For more information on this, see the \href{https://ai-bearing.github.io/Classes_and_Homework/After\%20College/Other\%20Notes/completely_positive_maps.pdf}{notes} on completely positive maps.
\subsection{Affiliated Operators}%
If $\left( M,\tau \right)$ is a tracial von Neumann algebra acting on a Hilbert space $\mathcal{H}$, we will show that the closed densely defined operators on $\mathcal{H}$ affiliated with $M$ behave ``nicely'' in a certain sense.

To start, we discuss some spectral theory of unbounded self-adjoint operators. If $x$ is a self-adjoint operator on a Hilbert space $\mathcal{H}$ --- that is, a densely defined possibly unbounded operator with $x = x^{\ast}$. The spectrum $\sigma\left( x \right)$ is a closed subset of $\R$. There is a $\ast$-homomorphism $f\mapsto f(x)$ taking $B_{\infty}\left( \sigma(x) \right)$ to $B\left( \mathcal{H} \right)$.

The functional calculus gives a spectral measure $E\colon \Omega\rightarrow E(\Omega)$, defined by $E(\Omega) = \chi_{\omega}(x)$, on the Borel subsets of $\sigma(x)$. Setting $E_t = E\left( \left( -\infty,t \right] \right)$, we have
\begin{align*}
  f(x) &= \int_{\sigma(x)}^{} f(t)\:dE_t.
\end{align*}
The functional calculus is extended to the algebra of all Borel functions $B\left(\sigma(x)\right)$ as follows. If $f\in B\left( \sigma(x) \right)$, then we have $\operatorname{dom}\left( f(x) \right)$ is given by
\begin{align*}
  \operatorname{dom}\left( f(x) \right) &= \set{\eta\in \mathcal{H} | \int_{\sigma(x)}^{} \left\vert f(t) \right\vert^2\:d \iprod{\eta}{E_t\eta} < \infty},
\end{align*}
and we define
\begin{align*}
  \iprod{\xi}{f(x)\eta} &= \int_{\sigma(x)}^{} f(t)\:d \iprod{\xi}{E_t\eta}.
\end{align*}
Whenever $f$ is real-valued, this is a closed densely defined operator, and we may write
\begin{align*}
  f(x) &= \int_{\sigma(x)}^{} f(t)\:dE_t\\
  x &= \int_{\sigma(x)}^{} t\:dE_t.
\end{align*}
Furthermore, we have
\begin{align*}
  \norm{f(x)\xi}^2 &= \int_{\sigma(x)}^{} \left\vert f(t) \right\vert^2\:d \iprod{\eta}{E_t\eta}.
\end{align*}
Now, if $y\in B\left( \mathcal{H} \right)$, then we say $y$ commutes with an unbounded operator $z$ if $\operatorname{dom}\left( yz \right)\subseteq \operatorname{dom}\left( zy \right)$ and $zy = yz$ on $\operatorname{dom}\left( yz \right)$. A bounded operator $y$ commutes with a self-adjoint operator $x$ if and only if it commutes with all its spectral projections. In that case, it commutes with $f(x)$ for every $f\in B\left( \sigma(x) \right)$.

There is an analogue of the polar decomposition for closed densely defined operators.
\begin{proposition}
  Let $x$ be a closed densely defined operator on $\mathcal{H}$. Then, the following hold:
  \begin{enumerate}[(i)]
    \item $x^{\ast}x$ is a positive self-adjoint operator;
    \item there exists a unique partial isometry $u$ such that $x = u\left\vert x \right\vert$ and $\ker\left( x \right) =\ ker\left( u \right)$, where $\left\vert x \right\vert = \left( x^{\ast}x \right)^{1/2}$.
  \end{enumerate}
\end{proposition}
We have $u^{\ast}u$ is the source projection onto $\ker\left( x \right)^{\perp}$ and $uu^{\ast}$ is the range projection onto $\img\left( x \right)$, denoted $s_r(x)$ and $s_{l}(x)$ respectively. We call them the right support and left support of $x$ respectively (note that they are Murray--von Neumann equivalent projections). An operator $y\in B(H)$ commutes with $x$ if and only if it commutes with both $u$ and $\left\vert x \right\vert$.
\begin{definition}
  Let $M$ be a von Neumann algebra on a Hilbert space $\mathcal{H}$. We say an (unbounded) operator $x$ is \textit{affiliated} with $M$ if for every unitary $u\in U\left( M' \right)$, we have $ux = xu$. We will write $x \widetilde{\in} M$
\end{definition}
That is, the operators $ux = xu$ have the same domain and coincide on this common domain. Denote by $ \widetilde{M} $ the set of all closed densely defined operators affiliated with $M$. The following is a consequence of the double commutant theorem and what we discussed earlier.
\begin{proposition}
  Let $M$ be a von Neumann algebra, $x$ a closed densely defined operator on $\mathcal{H}$, and $x = u\left\vert x \right\vert$ its polar decomposition. Then, $x\in \widetilde{M}$ if and only if $u$ and the spectral projections of $\left\vert x \right\vert$ are in $M$.
\end{proposition}
In particular, the left and right supports $s_r(x)$ and $s_l(x)$ are contained in $M$ whenever $x\in \widetilde{M}$.

In the case when $\tau$ is a trace, then $ \widetilde{M} $ has a particularly nice property.
\begin{proposition}
  Let $\left( M,\tau \right)$ be a tracial von Neumann algebra on $\mathcal{H}$, and let $x,y\in \mathcal{M}$ be such that $x\subseteq y$. Then, $x = y$.
\end{proposition}
Here, $x\subseteq y$ means that $\operatorname{dom}\left( x \right)\subseteq \operatorname{dom}\left( y \right)$ and $x = y|_{\operatorname{dom}\left( x \right)}$.
\begin{proof}
  Let
  \begin{align*}
    G(x) &= \set{\left( \xi,x\xi \right) | \xi\in \operatorname{dom}\left( x \right)}.
  \end{align*}
  Almost by definition, $x$ is a closed operator when $G(x)$ is a closed subspace of $\mathcal{H}\oplus \mathcal{H} = \mathcal{H}^{(2)}$. Similarly, we introduce $G(y)$ which is the graph of $y$. Let $\left[ G(x) \right]$ and $\left[ G(y) \right]$ be the projections of $\mathcal{H}^{(2)}$ onto $G(x)$ and $G(y)$ respectively.

  The algebra $\M_2\left( M \right)$ of $2\times 2$ matrices with entries in $M$ is a von Neumann subalgebra of $B\left( \mathcal{H}^{(2)} \right)$ with commutant given by
  \begin{align*}
    \M_2\left( M \right)' &= \set{ \operatorname{diag}\left( a,a \right) | a\in M' }.
  \end{align*}
  We claim that the projections $\left[ G(x) \right]$ and $\left[ G(y) \right]$ are in $\M_2\left( M \right)$. Since $ x \widetilde{\in} M $, we have for every $u'\in U\left( M' \right)$ that $u'x = xu'$, so that
  \begin{align*}
    \operatorname{diag}\left( u',u' \right) \left[ G(x) \right] &= \left[ G(x) \right] \operatorname{diag}\left( u',u' \right),
  \end{align*}
  so that $\left[ G(x) \right]\in \M_2\left( M \right)'' = \M_2\left( M \right)$, and similarly for $\left[ G(y) \right]$.

  Consider the projection
  \begin{align*}
    p_1 &= \begin{pmatrix}1 & 0 \\ 0 & 0\end{pmatrix}
  \end{align*}
  in $\M_2\left( M \right)$. Then, $p_1$ is the left support of the element $p_1\left[ G(x) \right]$ and $\left[ G(x) \right]$ is its right support. Similar reasoning holds for $p_1\left[ G(y) \right]$, meaning that $\left[ G(x) \right]\sim p_1\sim \left[ G(y) \right]$.

  Since $x\subseteq y$, we have $G(x)\subseteq G(y)$, so $\left[ G(x) \right]\leq \left[ G(y) \right]$. Since $\M_2(M)$ has a trace, it is finite, so $\left[ G(x) \right] = \left[ G(y) \right]$, meaning $G(x) = G(y)$, whence $x = y$.
\end{proof}
In general, if we consider two closed densely defined operators $x$ and $y$ on $\mathcal{H}$, the domains interact as
\begin{align*}
  \operatorname{dom}\left( x+y \right) &= \operatorname{dom}\left( x \right)\cap \operatorname{dom}\left( y \right),
\end{align*}
which in general is not dense in $\mathcal{H}$ (and can even be equal to $0$). Yet, if $x,y\in \widetilde{M}$, we will show that $x+y$ is densely defined and closable, which allows us to define an addition (and product) in $\widetilde{M}$.
\begin{lemma}
  Let $x\in \widetilde{M}$. Then, for every $\ve > 0$, there is a projection $p\in M$ such that $ p\mathcal{H}\subseteq \operatorname{dom}\left( x \right) $ and $\tau\left( 1-p \right)\leq \ve$.
\end{lemma}
\begin{proof}
  Let $x = u\left\vert x \right\vert$ be the polar decomposition of $x$, and set
  \begin{align*}
    p_n &= \chi_{[0,n]} \left( \left\vert x \right\vert \right).
  \end{align*}
  Then, $p_n\mathcal{H}\subseteq \operatorname{dom}\left( x \right)$, with $\lim_{n\rightarrow\infty}\tau\left( p_n \right) = 1$, so we may find some $n$ large enough with $1-\tau\left( p_n \right)\leq \ve$.
\end{proof}
\begin{lemma}
  Let $V$ be a subspace of $\mathcal{H}$ such that for every $\ve > 0$, there is a projection $p\in M$ with $p\mathcal{H}\subseteq V$ and $\tau\left( 1-p \right) \leq \ve$. Then, $V$ is dense in $\mathcal{H}$.
\end{lemma}
\begin{proof}
  It suffices to construct an increasing sequence $\left( q_n \right)_n$ of projections with $\sup_{n}q_n = 1$ and $q_n\mathcal{H}\subseteq V$.

  For each $k\geq 1$, find $p_k\in M$ with $p_k\mathcal{H}\subseteq V$ and $\tau\left( 1-p_k \right)\leq 2^{-k}$. Defining
  \begin{align*}
    q_n &= \inf_{k > n} p_k,
  \end{align*}
  we have $q_n\mathcal{H}\subseteq V$, and
  \begin{align*}
    \tau\left( 1-q_n \right) &= \tau\left( \sup_{k > n}\left( 1-p_k \right) \right)\\
                             &\leq \sum_{k > n} \tau\left( 1-p_k \right)\\
                             &\leq \sum_{k > n} 2^{-k}\\
                             &= 2^{-n},
  \end{align*}
  so since $\tau$ is normal, we have $\tau\left( 1-\sup_n q_n\right) = 0$, whence $\sup_{n}q_n = 1$.
\end{proof}
\begin{lemma}
  Let $\left( p_i \right)_{i\in I}$ be a family of projections in $\left( M,\tau \right)$. Then,
  \begin{align*}
    \tau\left( \sup_{i}p_i \right) &\leq \sum_{i}\tau\left( p_i \right).
  \end{align*}
\end{lemma}
\begin{proof}
  Using the Kaplansky formula, given two projections $p,q\in M$, we have
  \begin{align*}
    \tau\left( p\vee q-p \right) &= \tau\left( q-p\wedge q \right),
  \end{align*}
  whence $\tau\left( p\vee q \right)\leq \tau(p) + \tau(q)$. By induction, we have the inequality when $I$ is finite, and the case for infinite $I$ follows from the normality of $\tau$.
\end{proof}
\begin{theorem}
  Let $\left( M,\tau \right)$ be a tracial von Neumann algebra on $\mathcal{H}$.
  \begin{enumerate}[(i)]
    \item If $x\in \widetilde{M}$, then $x^{\ast}\in \widetilde{M}$.
    \item If $x,y\in \widetilde{M}$, then $x + y$ and $xy$ are densely defined and closable, with closures in $\widetilde{M}$.
    \item We have $\widetilde{M}$ equipped with these operations is a $\ast$-algebra.
  \end{enumerate}
\end{theorem}
\begin{proof}
  \begin{enumerate}[(i)]
    \item Since $x^{\ast}$ and $x$ have the same domain, this follows from taking adjoints in the definition of $x$.
    \item We start by showing that if $x,y\in \widetilde{M}$, then $x + y$ is densely defined. Given $\ve > 0$, let $p,q\in P(M)$ be such that $p\mathcal{H}\subseteq \operatorname{dom}\left( x \right)$, $q\mathcal{H}\subseteq \operatorname{dom}\left( y \right)$, and $\tau\left( 1-p \right)\leq \ve/2$, $\tau\left( 1-q \right)\leq \ve/2$. Then, we have
      \begin{align*}
        \left( p\wedge q \right)\mathcal{H} &= p\mathcal{H}\cap q\mathcal{H}\\
                                            &\subseteq \operatorname{dom}\left( x \right)\cap \operatorname{dom}\left( y \right)\\
                                            &= \operatorname{dom}\left( x + y \right),
      \end{align*}
      with
      \begin{align*}
        \tau\left( 1-p\wedge q \right) &= \tau\left( \left( 1-p \right)\vee \left( 1-q \right) \right)\\
                                       &\leq \tau\left( 1-p \right) + \tau\left( 1-q \right)\\
                                       &\leq \ve,
      \end{align*}
      so by one of the above lemmas, we have $x + y$ is densely defined, and is affiliated with $M$.
  \end{enumerate}
\end{proof}
\section{A Taste of Tomita--Takesaki Theory}%
The standard representation in fact generalizes to a wider setting --- namely, $\tau$ can be merely a faithful normal state instead of having to be a trace. This is the realm of Tomita--Takesaki Theory. We discuss a little bit of the ideas 

\nocite{takesaki_1,takesaki_2,von_neumann_algebras_peterson,ii1_factors_houdayer,ii1_factors_anantharaman_popa}
\printbibliography
\end{document}
